\section{Discussion}
The original hypothesis was that replacing fossil reformer heat with high-temperature nuclear heat while integrating CCS could materially improve SMR hydrogen performance in Singapore. The verified model shows that technical heat integration is plausible at screening level only after enforcing finite temperature approaches and consistent CCS duties. However, technical plausibility is not equivalent to assignment-level feasibility.

Across the evidence-backed uncertainty domain, the combined lifecycle and economic requirements are more restrictive than either engineering feasibility or stack capture alone. The 0/64 joint-pass result therefore shifts the interpretation from a presumed compliant nuclear solution to a conditional research scenario whose compliance is not demonstrated. This is a falsification outcome, not a failed research process.

The comparator evidence reinforces the need for common boundaries. Conventional Case 1A already clears the annual direct-abatement scale, but unmatched Singapore cost evidence prevents a defensible technology ranking. Future work should improve evidence where it changes decisions rather than broaden parameter ranges to search for a favourable nuclear point.
